% ID: SEL-D0016-BB-E04
% ESTADO: cotejada. Segunda pasada visual completa y borrador compilado revisados.
% Fuente: Selectividad/matematicas_aplicadas/matematicas_aplicadas_P_D.pdf
% SHA-256: 11865b83e509d9cc2723c752f7a81f714378e1b651f4a4d010dab4a4576b74c2
% Página: 1; bloque B; ejercicio original 4.

Se considera la función
\[
f(x)=\begin{cases}
ax+\dfrac{1}{2} & \text{si }x\leq -1\\
\dfrac{x+1}{x+3} & \text{si }-1<x\leq 1\\
x^2-bx & \text{si }x>1
\end{cases}
\]
\begin{enumerate}[label=\alph*)]
  \item \textbf{(1.5 puntos)} Halle $a$ y $b$ para que la función sea continua
  en todo su dominio. Para esos valores de $a$ y $b$,
  \textquestiondown{}es $f$ derivable en $x=-1$? \textquestiondown{}Y en $x=1$?
  \item \textbf{(0.5 puntos)} Para $a=-1$ y $b=4$, estudie la monotonía
  de la función $f$.
  \item \textbf{(0.5 puntos)} Para $a=-1$ y $b=4$, calcule
  $\displaystyle\int_1^2 f(x)\,dx$.
\end{enumerate}
